\section{Sticky priors for hidden Markov models}\label{sec:sticky}

Fox, Sudderth, Jordan and Willsky~\cite{fox2011} add a self-transition bias
$\kappa\ge0$ to the hierarchical Dirichlet process hidden Markov model
(HDP-HMM), so that the hidden chain stays in a state for long stretches.
With a finite vector $\beta$ of base weights on $d$ states, row $j$ of the
transition matrix has the Dirichlet law with parameters
$\omega\beta_1,\dots,\omega\beta_j+\kappa,\dots,\omega\beta_d$ in the
weak-limit approximation of~\cite{fox2011}. Here $\omega>0$ is the
concentration. By the neutrality of the Dirichlet law (write it as independent gamma variables divided by
their sum), the off-diagonal part of row $j$ divided by $1-\pi_{jj}$ is
independent of $\pi_{jj}$ and has the Dirichlet law with parameters
$(\omega\beta_l)_{l\ne j}$. So $\kappa$ only sets how sticky the chain is,
and the switching kernel is a random stochastic matrix with zero diagonal
whose law does not involve $\kappa$. We idealize by fixing the exit
probability of every state at $h=cL/N$, the scale of Kagey's crossing. In
the prior itself $1-\pi_{jj}$ has a Beta law and mixes all scales. That is a
different regime, and we do not treat it.

Throughout this section $\Pi=(\pi_{ij})$ is a stochastic matrix on $[d]$
with zero diagonal, $c>0$, and $h=cL/N$ with $0<h<1$. The transition matrix
$(1-h)I+h\Pi$ equals $I+hQ$ with $Q=\Pi-I$. So every exit rate is $q_i=1$,
$T=cL$ and $\tau=c$ (Example~\ref{ex:models}(iii)). We fix $\Pi$, which we
think of as one draw from the prior, and we assume that the start $\mu$ has
full support and that
\[
 \pi_{ij}>0\qquad\text{for all } i\ne j. \tag{A+}
\]
Under the prior, (A+) holds with probability $1$. For a face $F$ let $\Pi_F$
be the principal submatrix of $\Pi$ on $F$ and $\rho_F=\varrho(\Pi_F)$. Then
$Q_F=\Pi_F-I$, so $\lambda_F=1-\rho_F$, and $Q_F$ and $\Pi_F$ have the same
Perron vectors. A vertex has $\rho_F=0$, and $\rho_{[d]}=1$. If
$\lvert F\rvert\ge2$, then $\Pi_F$ is irreducible by (A+), and $\rho_F>0$
since $\rho_F$ is at least the smallest row sum of $\Pi_F$. If also
$F\ne[d]$, then $\rho_F<1$, because the Perron root of an irreducible
nonnegative matrix exceeds that of each proper principal
submatrix~\cite{bermanplemmons1994}. With $\tau=c$ the exponent of a face is
$\Phi_c(F)=c(1-\rho_F)+\lvert F\rvert-1$, and we put
\[
 c_-=\min_{\lvert F\rvert\ge2}\frac{\lvert F\rvert-1}{\rho_F},\qquad
 c_+=\max_{F\ne[d]}\frac{d-\lvert F\rvert}{1-\rho_F}.
\]
We call a global mode of $P_N$ a \emph{MAP composition}. With no data the
posterior is the prior, so a MAP composition is the most likely value of the
occupation counts under the prior.

\begin{corollary}[support and location of the MAP]\label{cor:map}
Assume \textup{(A+)}, let $\mu$ have full support, and fix $c>0$. Let
$\mathcal F^*_c$ be the set of faces that minimize $\Phi_c$. There is
$N_0(c)$ such that for $N\ge N_0(c)$ every MAP composition $n^*$ satisfies
$\supp n^*\in\mathcal F^*_c$ and $n^*/N=\alpha^*_F+O(1/L)$ with
$F=\supp n^*$. If $\mathcal F^*_c=\{F^*\}$, then
$n^*/N\to\alpha^*_{F^*}=l\circ r$, where $l$ and $r$ are the left and right
Perron vectors of $\Pi_{F^*}$ with $l\cdot r=1$. No reversibility is needed.
\end{corollary}

\begin{proof}
Every face is irreducible, since $G_F$ is complete by (A+) when
$\lvert F\rvert\ge2$. Every face is also admissible, since a path that
visits the states of $F$ one after another has positive probability when
$\mu$ has full support and $0<h<1$. So
Theorem~\ref{thm:first}(iii) with $\tau=c$ gives $\supp n^*\in\mathcal F^*_c$
for large $N$. Its assumption $hq_{\max}\le\frac12$ holds for large $N$. The
faces of $\mathcal F^*_c$ are irreducible (as Theorem~\ref{thm:hull}(c) also
says), so Theorem~\ref{thm:first}(iv) gives $n^*/N=\alpha^*_F+O(1/L)$. For
$\lvert F\rvert\ge2$ the Perron vectors of $Q_F$ are those of $\Pi_F$, so
$\alpha^*_F=l\circ r$. For a vertex $\{i\}$, $n^*=Ne_i$. If
$\mathcal F^*_c=\{F^*\}$, then $\supp n^*=F^*$ and the error tends to $0$.
\end{proof}

The corollary is pointwise in $c$. At $c=c_-$ the vertices tie with the faces
that attain the minimum in $c_-$. Every face is irreducible, so
Theorem~\ref{thm:window} applies. Part (iii) of that theorem, with a
fixed large $t=K$, shows that for large $N$ the MAP is not a single state at
$c=c_-\bigl(1-\frac{\log L-2K}{2L}\bigr)$, which is below $c_-$. So
$N_0(c)\to\infty$ as $c\uparrow c_-$. Theorem~\ref{thm:window} also describes
the window at $c_+$, with the coefficient $\frac12$ of $\log\log N$.

\begin{proposition}[thresholds and the direct jump]\label{prop:thresholds}
Assume \textup{(A+)} and $d\ge3$.
\begin{enumerate}[(a)]
\item If $0<c<c'$, $F$ minimizes $\Phi_c$ and $G$ minimizes $\Phi_{c'}$,
then $\lvert F\rvert\le\lvert G\rvert$ and $\lambda_F\ge\lambda_G$.
\item The minimizers of $\Phi_c$ are exactly the vertices if and only if
$c<c_-$, and no vertex minimizes $\Phi_c$ if $c>c_-$. The face $[d]$ is the
unique minimizer if and only if $c>c_+$, and $[d]$ does not minimize $\Phi_c$
if $c<c_+$.
\item $1<c_-\le d-1\le c_+$.
\item $c_-=d-1$ if and only if $c_+=d-1$, and both hold if and only if
\begin{equation}\label{eq:directjump}
 \rho_F\le\frac{\lvert F\rvert-1}{d-1}\qquad\text{for every face } F.
\end{equation}
If \eqref{eq:directjump} fails, then $c_-<d-1<c_+$, and for
$c\in(c_-,c_+)$ every minimizer $F$ of $\Phi_c$ has
$2\le\lvert F\rvert\le d-1$.
\end{enumerate}
Consequently, let $\mu$ have full support, fix $c$ and let $N\ge N_0(c)$. If
$c<c_-$, the MAP compositions are the $Ne_i$ with $\mu_i=\max_l\mu_l$. If
$c>c_+$, every MAP composition has full support. If $c_-<c<c_+$, every MAP
composition has a support $F$ with $2\le\lvert F\rvert\le d-1$, and $F=F^*$
if $F^*$ is the unique minimizer of $\Phi_c$.
\end{proposition}

\begin{proof}
(a) Every face is admissible, so this is Theorem~\ref{thm:hull}(b).

(b) A vertex has $\Phi_c=c$. For $\lvert F\rvert\ge2$,
$\Phi_c(F)-c=\rho_F\bigl((\lvert F\rvert-1)/\rho_F-c\bigr)$ with
$\rho_F>0$. So if $c<c_-$, every face of size at least $2$ is worse than the
vertices, which all tie. The face attaining the minimum in $c_-$ ties with
the vertices at $c=c_-$ and beats them for $c>c_-$. In the same way, for
$F\ne[d]$ we have $\Phi_c(F)-(d-1)=(1-\rho_F)\bigl(c-(d-\lvert F\rvert)/
(1-\rho_F)\bigr)$ with $1-\rho_F>0$, and $\Phi_c([d])=d-1$. This gives the
claims about $[d]$, using the face that attains the maximum in $c_+$.

(c) If $2\le\lvert F\rvert\le d-1$, then $\rho_F<1$ gives
$(\lvert F\rvert-1)/\rho_F>\lvert F\rvert-1\ge1$, and $F=[d]$ gives
$d-1\ge2$. So $c_->1$. The face $[d]$ gives $c_-\le d-1$, and a vertex gives
$c_+\ge d-1$.

(d) By (c), $c_+=d-1$ if and only if $(d-\lvert F\rvert)/(1-\rho_F)\le d-1$
for all $F\ne[d]$, and $c_-=d-1$ if and only if
$(\lvert F\rvert-1)/\rho_F\ge d-1$ for all $F$ with $\lvert F\rvert\ge2$.
Both inequalities say $(d-1)\rho_F\le\lvert F\rvert-1$, which holds for
every vertex ($\rho_F=0$) and for $[d]$ ($\rho_F=1$). This proves the
equivalences. If \eqref{eq:directjump} fails, (c) gives $c_-<d-1<c_+$, and
for $c\in(c_-,c_+)$ part (b) excludes the vertices and $[d]$.

The consequences follow from Corollary~\ref{cor:map}, since
$P_N(Ne_i)=\mu_i(1-h)^{N-1}$.
\end{proof}

We say that $\Pi$ has a \emph{direct jump} when \eqref{eq:directjump} holds.
Then at first order the MAP goes from one state straight to all $d$ states
at $c=d-1$. The chance that $m$ switches in a row from inside $F$ all land in
$F$ is of order $\rho_F^m$. For the uniform kernel
$\pi_{ij}=1/(d-1)$ every row sum of $\Pi_F$, and hence $\rho_F$, equals
$(\lvert F\rvert-1)/(d-1)$. So \eqref{eq:directjump} says that no set of
states is more closed than under the uniform kernel. For that kernel all
faces tie at $c=d-1$, which is the degeneracy of
Proposition~\ref{prop:complete}(a), since $Q=(J-dI)/(d-1)$ is the generator
of Paper~B~\cite{paperB} divided by $d-1$ (Example~\ref{ex:models}(ii)).

\begin{proposition}[when the direct jump happens]\label{prop:directjump}
\begin{enumerate}[(a)]
\item Let $d\ge3$, assume \textup{(A+)}, and let $\Pi$ be symmetric. Then
$\Pi$ has a direct jump if and only if $\pi_{ij}=1/(d-1)$ for all $i\ne j$.
\item Let $d=3$ and assume \textup{(A+)}. Then $\Pi$ has a direct jump if and
only if $\pi_{ij}\pi_{ji}\le\frac14$ for all $3$ pairs. For example, the
cyclic kernel $\pi_{i,i+1}=a$, $\pi_{i,i-1}=1-a$ (indices mod $3$) has a
direct jump for every $0<a<1$.
\item Let $d=3$ and let the base weights be uniform. In the weak-limit sticky
HDP-HMM the rows of $\Pi$ are then independent, and the first off-diagonal
entry of each row has the law $\mathrm{Beta}(\omega/3,\omega/3)$. Write
$\Prob_\omega$ for this law of $\Pi$ and $D$ for the event that $\Pi$ has a
direct jump. Then $\Prob_\omega(D)\to\frac14$ as $\omega\to0$.
\item In the setting of \textup{(c)},
$\Prob_\omega(D)=\frac{3\sqrt3}{8\pi\omega}\bigl(1+o(1)\bigr)$ as
$\omega\to\infty$.
\end{enumerate}
\end{proposition}

\begin{proof}
For a pair $\{i,j\}$ the eigenvalues of $\Pi_{\{i,j\}}$ are
$\pm\sqrt{\pi_{ij}\pi_{ji}}$, so $\rho_{\{i,j\}}=\sqrt{\pi_{ij}\pi_{ji}}$.

(a) Here $\rho_{\{i,j\}}=\pi_{ij}$. The rows sum to $1$, so
$\sum_{i<j}\pi_{ij}=d/2$, and the average of $\pi_{ij}$ over the $\binom d2$
pairs is $1/(d-1)$. Unless all $\pi_{ij}$ are equal, some pair has
$\rho_{\{i,j\}}>1/(d-1)$, and \eqref{eq:directjump} fails. The uniform kernel
satisfies \eqref{eq:directjump} with equality, as shown above.

(b) For $d=3$ the condition \eqref{eq:directjump} is automatic for vertices
and for $[3]$, and for a pair it reads $\pi_{ij}\pi_{ji}\le\frac14$. The
cyclic kernel satisfies (A+) since $0<a<1$, and every pair has
$\pi_{i,i+1}\pi_{i+1,i}=a(1-a)\le\frac14$.

(c) The law of $\Pi$ comes from the neutrality property above. Let
$B_1,B_2,B_3$ be independent with law $\mathrm{Beta}(\omega/3,\omega/3)$ and
put $\pi_{12}=B_1$, $\pi_{13}=1-B_1$, $\pi_{21}=B_2$, $\pi_{23}=1-B_2$,
$\pi_{31}=B_3$ and $\pi_{32}=1-B_3$. Then (A+) holds with probability $1$,
and by (b) the event $D$ is
$\{B_1B_2\le\frac14,\ (1-B_1)B_3\le\frac14,\ (1-B_2)(1-B_3)\le\frac14\}$. As
$\omega\to0$ each $B_j$ converges in law to a fair coin on $\{0,1\}$. Under
the limit law the 3 products take only the values $0$ and $1$, so the
boundary of $D$ has limit probability $0$, and $\Prob_\omega(D)$ converges
to the limit probability of $D$. Of the $8$ points of $\{0,1\}^3$, only
$(1,0,1)$ and $(0,1,0)$ make all 3 products $0$. So the limit is $\frac14$.

(d) This is proved in Appendix~\ref{app:hdp}.
\end{proof}

The 2 points in (c) are the directed $3$-cycles, the only deterministic
kernels in which no pair of states traps the chain. By (a) and (d), the tie
of all faces that comes from the uniform kernel is not typical. Numerically
$\Prob_\omega(D)=0.115$ at $\omega=1$ (Monte Carlo with $10^6$ draws). So with
probability about $0.88$ a random kernel has a window of $c$ in which the MAP
sits on a pair of states (Proposition~\ref{prop:thresholds}(d)).
Figure~\ref{fig:sticky} shows $\Prob_\omega(D)$ against $\omega$. Quadrature
values suggest $\Prob_\omega(D)=\frac{3\sqrt3}{8\pi\omega}\bigl(1-\frac9
{8\omega}+O(\omega^{-2})\bigr)$, but we have not proved the second term.

\begin{figure}[t]
\centering
\includegraphics{fig_sticky}
\caption{The probability $\Prob_\omega(D)$ that a random sticky kernel on $3$
states with uniform base weights has a direct jump
(Proposition~\ref{prop:directjump}), against the concentration $\omega$. The
open circles are Monte Carlo estimates with $10^6$ draws, and their error bars
of one standard error are smaller than the circles. The filled squares are
quadrature values for $\omega\ge100$. The dotted line is the limit $\frac14$ as
$\omega\to0$, and the dashed line is the asymptote $3\sqrt3/(8\pi\omega)$ of
Proposition~\ref{prop:directjump}(d).}
\label{fig:sticky}
\end{figure}

When the first switch of the MAP goes from one state to a symmetric pair, it
has an exact description at every $N$. We use $S_{a,b}$, $S_N$, $u_N$, $p_N$
and $c_0$ from Section~\ref{sec:papersAB}.

\begin{proposition}[an exact link to Kagey's crossing]\label{prop:sticky-pair}
Let $i\ne j$ satisfy $\pi_{ij}=\pi_{ji}=\pi_0>0$ and $\mu_i=\mu_j$. We do not
assume \textup{(A+)}. For $1\le m\le N-1$,
\[
 \frac{P_N(me_i+(N-m)e_j)}{P_N(Ne_i)}=S_{m,N-m}(u'),\qquad
 u'=\frac{h\pi_0}{1-h}.
\]
\begin{enumerate}[(a)]
\item Among the compositions with support exactly $\{i,j\}$, the most likely
ones are the balanced ones, with $m=\lfloor N/2\rfloor$ or
$m=\lceil N/2\rceil$.
\item The balanced composition $\lfloor N/2\rfloor e_i+\lceil N/2\rceil e_j$
is more likely than $Ne_i$ if and only if $u'>u_N$, that is, if and only if
$c>c^{\rm pair}(N):=Nu_N/((\pi_0+u_N)L)$.
\item If also $\mu_i=\mu_j=\max_l\mu_l$, then for every $N\ge2$ and every
$c$ with $c^{\rm pair}(N)<c<N/L$, no MAP composition is a single state.
\end{enumerate}
Moreover,
\[
 c^{\rm pair}(N)=\frac1{\pi_0}\Bigl[1-\frac{\frac12\log L-c_0}{L}
 +O\Bigl(\frac{\log L}{L^2}\Bigr)\Bigr],
\]
which tends to the first-order value $1/\rho_{\{i,j\}}=1/\pi_0$.
\end{proposition}

\begin{proof}
Every exit rate is $1$, so the identity is Proposition~\ref{prop:pair} with
$a=\pi_0$ and $q=1$, whose condition $qh<1$ is $h<1$. Note that
$P_N(Ne_i)=\mu_i(1-h)^{N-1}>0$ and $u'>0$.

(a) By \cite[Proposition~6.1]{paperA}, if $a<b-1$, then $S_{a+1,b-1}$
dominates $S_{a,b}$ coefficientwise with at least one strict coefficient.
So $S_{m+1,N-m-1}(u')>S_{m,N-m}(u')$ whenever $2m<N-1$, and $S_{m,N-m}(u')$
strictly increases in $m$ up to $m=\lfloor N/2\rfloor$. The run formula is
symmetric, so $S_{a,b}=S_{b,a}$, and $S_{m,N-m}(u')$ strictly decreases from
$m=\lceil N/2\rceil$ on.

(b) For $m=\lfloor N/2\rfloor$ the ratio is $S_N(u')$. By
\cite[Theorem~5.1]{paperA} and its proof, $S_N$ is strictly increasing on
$[0,\infty)$ with the unique root $u_N$ of $S_N(u)=1$. So the balanced
composition wins if and only if $u'>u_N$, that is $h\pi_0>u_N(1-h)$, which
is $c>c^{\rm pair}(N)$.

(c) Every vertex has $P_N(Ne_l)=\mu_l(1-h)^{N-1}\le P_N(Ne_i)$, and by (b)
$P_N(Ne_i)$ is less than the probability of the balanced composition. The
range of $c$ is not empty since $u_N/(\pi_0+u_N)<1$.

For the expansion, \cite[Theorem~5.2]{paperA} gives
$N(1-p_N)=L-\frac12\log L+c_0+O(\log L/L)$. Since $u_N=(1-p_N)/p_N$ and
$1-p_N=O(L/N)$, also $Nu_N=L-\frac12\log L+c_0+O(\log L/L)$ and
$u_N=O(L/N)$. Then
$c^{\rm pair}(N)=\frac{Nu_N}{\pi_0L}(1+u_N/\pi_0)^{-1}$ gives the claim,
since $L/N=O(\log L/L^2)$.
\end{proof}

So for a sticky prior, one state against one pair with $\pi_{ij}=\pi_{ji}$
is exactly Kagey's crossing, at every $N$. With a uniform start, part
(c) says that the MAP stops being a single state no later than
$c^{\rm pair}(N)$.

The last result says where the prior mass is. Put $\Lambda_F=1-h\lambda_F$.

\begin{proposition}[face masses]\label{prop:mass}
Assume \textup{(A+)} and $0<h<1$. For every face $F$ with
$\lvert F\rvert\ge2$,
\[
 \frac{\min r_F}{\max r_F}\,\mu(F)\,\Lambda_F^{N-1}
 \le\Prob(\supp n\subseteq F)\le
 \frac{\max r_F}{\min r_F}\,\mu(F)\,\Lambda_F^{N-1},
\]
and for fixed $c$, $\Lambda_F^{N-1}=N^{-c\lambda_F}e^{O(L^2/N)}$. Also
$\Prob(\text{one state})=(1-h)^{N-1}\sim N^{-c}$ and
$\Prob(\supp n=[d])\to1$.
\end{proposition}

\begin{proof}
Let $U_F=(1-h)I+h\Pi_F$, which is nonnegative since $h<1$. Then
$\Prob(\supp n\subseteq F)=\mu_F^{\top}U_F^{N-1}\ones$, where $\mu_F$ is the
restriction of $\mu$ to $F$. The Perron vector $r_F>0$ of $\Pi_F$ satisfies
$U_Fr_F=\Lambda_Fr_F$. Entrywise $r_F/\max r_F\le\ones\le r_F/\min r_F$, and
$U_F^{N-1}$ preserves these inequalities. So $U_F^{N-1}\ones$ lies between
$\Lambda_F^{N-1}r_F/\max r_F$ and $\Lambda_F^{N-1}r_F/\min r_F$, and pairing
with $\mu_F$ gives the bounds. Next,
$(N-1)\log(1-h\lambda_F)=-(N-1)h\lambda_F+O(Nh^2)=-c\lambda_FL+O(L^2/N)$.
A vertex composition is a single path, so
$\Prob(\text{one state})=\sum_i\mu_i(1-h)^{N-1}=(1-h)^{N-1}$. Finally, every
$F\ne[d]$ has $\lambda_F>0$, so $\Prob(\supp n\ne[d])$ is at most a finite
sum of terms $O(N^{-c\lambda_F})$.
\end{proof}

\begin{remark}[MAP versus mass]\label{rem:mapmass}
Let $d\ge3$. For fixed $c<c_+$ the face $[d]$ does not minimize $\Phi_c$, so
for large $N$ every MAP composition has a proper support, while
$\Prob(\supp n=[d])\to1$. For $c<c_-$ the most likely composition uses one
state, and yet that event has probability about $N^{-c}$. The reason is that
a composition on a face of size $k$ is one of about $N^{k-1}$ lattice points.
So a face can carry most of the mass while each of its compositions is less
likely than a vertex. Table~\ref{tab:sticky} shows this for the kernel
$\Pi_1$ with rows $(0,0.9,0.1)$, $(0.9,0,0.1)$, $(0.5,0.5,0)$ and the
uniform start. Its pair $\{1,2\}$ is a well-separated community with
$\rho_{\{1,2\}}=0.9$. In part (a) the first switch of the MAP agrees with
$c^{\rm pair}(N)$ to the digits shown. In part (b), at $c=0.8$, the MAP is a
single state (all 3 vertices tie), yet $\Prob(\text{one state})=0.0122$. At
$c=8$ the first order puts the MAP on $\{1,2\}$, but at $N=240$ the second
switch has already happened. So the MAP is a poor summary of a sticky prior,
and face selection says when it misleads.
\end{remark}

\begin{table}[t]
\centering\small
\begin{tabular}{@{}lcccc@{}}
\multicolumn{5}{@{}l}{(a) Switches of the MAP support along $c$}\\
\toprule
$N$ & 60 & 120 & 240 & 480\\
\midrule
first switch, vertex to $\{1,2\}$ & 0.8248 & 0.8477 & 0.8673 & 0.8838\\
$c^{\rm pair}(N)$ & 0.8248 & 0.8477 & 0.8673 & 0.8838\\
second switch, $\{1,2\}$ to $[3]$ & 4.804 & 5.340 & 5.779 & 6.139\\
\bottomrule
\end{tabular}

\medskip
\begin{tabular}{@{}rlcccc@{}}
\multicolumn{6}{@{}l}{(b) The MAP and the prior mass at $N=240$}\\
\toprule
$c$ & MAP composition & one state & $\{1,2\}$ & other pairs & $[3]$\\
\midrule
0.5 & one state (3-way tie) & 0.0643 & 0.465 & 0.092 & 0.379\\
0.8 & one state (3-way tie) & 0.0122 & 0.423 & 0.034 & 0.531\\
3 & $(120,120,0)$ & $4.3\times10^{-8}$ & 0.129 & $3.3\times10^{-6}$ & 0.871\\
8 & $(111,111,18)$ & $1.2\times10^{-21}$ & 0.0081 & $2.1\times10^{-16}$ &
0.992\\
\bottomrule
\end{tabular}
\caption{The sticky chain with kernel $\Pi_1$ and the uniform start, for
which $c_-=10/9$ and $c_+=10$. (a) The 2 switches of the MAP support along
$c$, from exact dynamic programming, and the exact value $c^{\rm pair}(N)$ of
Proposition~\ref{prop:sticky-pair}. In these computations the faces
$\{1,3\}$ and $\{2,3\}$ never hold the MAP. (b) At $N=240$, the MAP
composition and the prior probabilities of the possible supports. The column
for the other pairs gives the total for $\{1,3\}$ and $\{2,3\}$.}
\label{tab:sticky}
\end{table}
